\documentclass{article}
\usepackage{amsmath, amssymb, amsthm}
\usepackage{parskip}
\usepackage{setspace}
\usepackage{titlesec}
\onehalfspacing
\setlength{\jot}{10pt}
\titleformat{\section}{\normalsize\bfseries}{\thesection}{1em}{}
\titleformat{\subsection}{\small\bfseries}{\thesubsection}{1em}{}

\newtheorem{definition}{Definition}
\theoremstyle{remark}
\newtheorem*{remark}{Remark}

\title{The Gradient of the Logistic Function}
\author{Carter Deslauriers}
\date{\today}

\begin{document}
\maketitle

\section{What does a Logistic Function do?}

% SET UP
\subsection{Setup}

Take one row of observations. Let there be $d$ observations; where each $x_j$
is the z-score for jth observation (the standard deviations from the mean), 

\begin{equation*}
    x_1, x_2, \ldots , x_d
\end{equation*} 

and a single classification observation $y$ where

\begin{equation*}
    y \in \{0, 1\}
\end{equation*} 

You can think of $y$ as yes or no, or true or false.

% MOTIVATION FOR Z
\subsection{Motivation for z}

Our goal is simply take these $x_j$ and predict $y$ as best as possible.
So intuitvely we can take a single observation call it $x_2$. Choose an arbitray break
$x_2 > 1$ say. Let our prediction by $1$ for $y$ when $x_2 > 1$.
But this doens't take into account all our observations. We want to take
into consideration all of our variables when making a desicion. So let
z be the summation of the $d$ observations.

\begin{equation*}
z = x_1 + x_2 +, \ldots ,+ x_d + b
\end{equation*} 

We introduced a b
which is a bias for the guess we want. You can imagine b as, if all $x_j$ are 0 then
what would we like our guess be. 

% MOTIVATION FOR WEIGHTS
\subsection{Motivation for weights}

Now it might be the case that some observations are more important
than others, or interact with y in different ways. For example your length of credit
history might be more imporant than your hair color in predicting if you will
repay or not. 

So assign a weight $w_j$ to each observation $x_j$. So we have now

\begin{equation}
    z = w_1x_1 + w_2x_2 + ... + w_dx_d + b
\end{equation}

But how do we get these weights? We need to conlude based on the current weights
how confident we are that $y = 1$ ie given our current value of z how confident are we
that this observation is a true $y = 1$. Well z is arbitary, $w_1$ can be as big as we want,
z can be $100$, $-0.5$ so on; what does that mean?

% SIGMA(Z)
\subsection{$\sigma(z)$}

Take a function 

\begin{equation}
    \sigma(z) = \frac{1}{1 + e^{-z}}
\end{equation}

and let this be the predicted probability that $y = 1$. 
Specifcally let 

\begin{equation}
    p = \sigma(z) = \frac{1}{1 + e^{-z}}
\end{equation}

To get a concept of a probability we must have the function strictly $p \in [0, 1]$

% PROOF FOR P
\subsection{Proof $p \in (0, 1)$}

Observe that for every $\epsilon > 0$ theres exists $N \in \mathbb{N}$ such that whenever
$z \geq N$ it follows

\begin{equation*}
    \left|\frac{1}{1 + e^{-z}} - 1\right| < \epsilon
\end{equation*}

specifcally choose

\begin{equation*}
    N > \log\!\left(\dfrac{1}{\epsilon}\right)
\end{equation*}

To see this

\begin{align*}
\left|\frac{1}{1 + e^{-z}} - 1\right|
  &= \left|\frac{1}{1 + e^{-z}} - \frac{1 + e^{-z}}{1 + e^{-z}}\right| \\
  &= \left|\frac{-e^{-z}}{1 + e^{-z}}\right| \\
  &= \frac{e^{-z}}{1 + e^{-z}}
\end{align*}

Where the last line follows since $e^{x} > 0$ always.

This gives

\begin{equation*}
    {1 + e^{-z}} > 1
\end{equation*}

giving 

\begin{equation*}
    \frac{e^{-z}}{1 + e^{-z}} < e^{-z} 
\end{equation*}

Now 

\begin{align*}
e^{-z} < \epsilon
  &\iff -z < \log(\epsilon) \\
  &\iff z > -\log(\epsilon) \\
  &\iff z > \log(\frac{1}{\epsilon})
\end{align*}

Meaing whenver $z \geq N$ where $N > \log(\frac{1}{\epsilon})$ we have 

\begin{equation*}
\left|\frac{1}{1 + e^{-z}} - 1\right| < \epsilon 
\end{equation*}

giving 
\begin{equation*}
    p \to 1 \text{ as } z \to \infty
\end{equation*}

For the other direction observe that for every $\epsilon > 0$ 
theres exists $N \in \mathbb{N}$ such that whenever$z \leq -N$ it follows

\begin{equation*}
\left|\frac{1}{1 + e^{-z}} - 0\right| < \epsilon 
\end{equation*}

specifcally choose $N > -\log(\frac{1 - \epsilon}{\epsilon})$

To see this

\begin{equation*}
\left|\frac{1}{1 + e^{-z}} - 0\right| = \frac{1}{1 + e^{-z}} 
\end{equation*}

Which follows from $e^{x} > 0$ always. Again giving

\begin{equation*}
    1 < \epsilon({1 + e^{-z}})
\end{equation*}

and 

\begin{equation*}
    \frac{1 - \epsilon}{\epsilon} < e^{-z}
\end{equation*}

taking log of both sides and multiplying by $-1$

\begin{equation*}
\log(\frac{1 - \epsilon}{\epsilon}) < -z \iff -\log(\frac{1 - \epsilon}{\epsilon}) > z
\end{equation*} 

Meaning whenver $z \leq -N$ where $N > -\log(\frac{1 - \epsilon}{\epsilon})$ we have 

\begin{equation*}
\left|\frac{1}{1 + e^{-z}} - 0\right| < \epsilon
\end{equation*}

giving

\begin{equation*}
p \to 0 \text{ as } z \to -\infty 
\end{equation*}


Now that we have proved $p \in (0, 1)$ and, we have defined it as the probabilty
that $y = 1$ for our model, we can define the loss that we are looking to minimize.

% LOSS
\subsection{Loss}

Intuitively we would want the distance from our predicted probability
to the true value of $y=0$ or $y=1$ But an absoulte difference 
doens't account for the confidence difference between a prediction of .99 
and .999, which are 0.01 and 0.001 away from 1 but are dramatically
different in confidence. Therefore we define the loss using logs. Specifically we
let the loss be

\begin{equation*}
L = -\log(p)
\end{equation*}

if $y=1$ and 

\begin{equation*}
L -\log(1 - p)
\end{equation*}
if $y=0$. 

In closed form our loss looks like 

\begin{equation}
    L = -[y\log(p) + (1 - y)\log(1-p)]
\end{equation}

% GRADIENT
\subsection{Gradient}

As stated before we want to minimze our loss. So how can we do this. 
We want to tune our weights as optimally as possible. Therefore
we use the gradient specifcally 

\begin{equation*}
\frac{\partial L}{\partial w_j}
\end{equation*}

But we defined L in terms of $y$ and $p$ 
so we must chain through to get our derivative.

Observe by the chain rule 

\begin{equation*}
\frac{\partial L}{\partial w_j} = 
\frac{\partial L}{\partial p}\frac{\partial p}{\partial z}\frac{\partial z}{\partial w_j}
\end{equation*}

Moving from right to left we have 

\begin{equation*}
\frac{\partial z}{\partial w_j} = x_j
\end{equation*}

since taking the partial derivative leaves all other weights as constants

Next 

\begin{align*}
\frac{\partial p}{\partial z}
  &= \sigma'(z) \\
  &= -1(1 + e^{-z})^-2e^{-z} \\
  &= \frac{e^{-z}}{(1 + e^{-z})^2} 
\end{align*}
  
Observe that 

\begin{equation*}
\frac{e^{-z}}{(1 + e^{-z})^2} = \frac{e^{-z}}{(1 + e^{-z})}\frac{1}{{(1 + e^{-z})}}
\end{equation*}

where the right side is again
$\sigma(z)$ and the left side is $1 - \sigma(z)$ since 

\begin{align*}
1 - \sigma(z)
  &= 1 - \frac{1}{1 + e^{-z}}
  &= \frac{1 + e^{-z}}{1 + e^{-z}} - \frac{1}{1 + e^{-z}}
  &= \frac{e^{-z}}{(1 + e^{-z})} 
\end{align*}
  
Giving 

\begin{equation*}
p' = p(1-p)
\end{equation*}

Finally 

\begin{equation*}
\frac{\partial L}{\partial p} = \frac{-y}{p} + \frac{1 - y}{1 - p}
\end{equation*}

Now chaining all three together gives 

\begin{align*}
(\frac{-y}{p} + \frac{1 - y}{1 - p})p(1-p)x_j
  &= (-y(1-p) + (1-y)p)x_j 
  &= (-y + py + p -py)x_j 
  &= (p - y)x_j
\end{align*}

giving the gradient as

\begin{equation}
(p - y)x_j
\end{equation}

For b simply observe

\begin{equation*}
\frac{\partial z}{\partial b} = 1
\end{equation*}

and therefore the gradient is 
\begin{equation}
(p - y)
\end{equation}

We now have the optimal unit of movement of each $w_j$ in order to most efficencly
decrease our loss. 

% N ROWS OF DATA
\subsection{N rows of Data}

In general we have $d$ columns of observations and $n$ rows of data. 
We can then simply think of each vairbale described above as a vector. 
Where w is the weight vecotr in which we add b.
p is the vector of probabilites one for each row $p_i$.
y is the vector of observations for rows i where $y_i$ is the value 0 or 1
descirbed throughout.
To get our final gradient we take Loss as the avergae loss of the rows.

\begin{equation}
L = \frac{1}{n}\sum_{i=1}^{n}L_i
\end{equation}

where 

\begin{equation*}
L_i = (p_i - y_i)x_ij
\end{equation*}

for each observation j. Then to find the partial derivavite of
the gradient we have 

\begin{equation*}
\frac{\partial L}{\partial w_j} =
\frac{1}{n}\sum_{i=1}^{n}\frac{\partial L_i}{\partial w_ij}
\end{equation*}

since the partial deriative of the mean is the mean of the partial derivatives.

% OPTIMUM
\subsection{Optimum}

In general at optimum we have 

\begin{equation}
\sum_{i=1}^{n}p_i = \sum_{i=1}^{n}y_i
\end{equation}

Which follows from

\begin{equation*}
L = \frac{1}{n}\sum_{i=1}^{n}L_i
\end{equation*}

where

\begin{equation*}
L_i = (p_i - y_i)x_ij
\end{equation*}

and for b

\begin{equation*}
L_i = (p_i - y_i)
\end{equation*}

giving

\begin{align*}
\frac{1}{n}\sum_{i=1}^{n} (p_i - y_i) = 0
  &\iff \frac{1}{n}\sum_{i=1}^{n} p_i - \frac{1}{n}\sum_{i=1}^{n} y_i = 0 \\
  &\iff \frac{1}{n}(\sum_{i=1}^{n} p_i = \frac{1}{n}\sum_{i=1}^{n} y_i) \\
  &\iff \sum_{i=1}^{n} p_i = \sum_{i=1}^{n} y_i
\end{align*}


Where our logistic regression will produce its most optimal predictions
based on our assumptions and data input.

We can then use the model to predict further classifcations by training on data
we believe to be representative of the further data we are trying to predict.

\end{document}